\documentclass[11pt,a4paper]{article}
\usepackage[T1]{fontenc}
\usepackage{lmodern,amsmath,amssymb,amsthm,mathtools,microtype,booktabs}
\usepackage[margin=26mm]{geometry}
\usepackage[colorlinks=true,linkcolor=blue,citecolor=blue,urlcolor=blue]{hyperref}
\setlength{\emergencystretch}{2em}
\newtheorem{theorem}{Theorem}[section]
\newtheorem{proposition}[theorem]{Proposition}
\newtheorem{lemma}[theorem]{Lemma}
\newtheorem{corollary}[theorem]{Corollary}
\theoremstyle{definition}\newtheorem{definition}[theorem]{Definition}
\theoremstyle{remark}\newtheorem{remark}[theorem]{Remark}
\newcommand{\Q}{\mathbb Q}\newcommand{\Z}{\mathbb Z}
\newcommand{\A}{\mathbb A}\newcommand{\sqcl}[1]{[#1]}
\DeclareMathOperator{\Tr}{Tr}\DeclareMathOperator{\N}{N}
\title{Joint realization of compatible pole conditions}
\author{}\date{14 November 2026}
\begin{document}\maketitle
\begin{abstract}
An arbitrary finite collection of pole inputs with matching input and
complement square classes can use a single pair of retained rational
parameters. We prove explicit invertible changes of pole witnesses for
both pole systems used in this project. Additional square-class tests are
finite algebraic choices, not new retained parameters. The older system
requires matching classes of $b$ and $17-b$; the congruent-number system
requires matching classes of $b$ and $1-b$. The former uses only the
previously retained pole-realization input. For the latter we state the
exact use of Jaswal's simultaneous-twist theorem and prove a more general
conditional common-parameter theorem.

These are reductions of pole subsystems, not new records for defining
$\Z$ in $\Q$. In the present thirteen-variable completeness construction,
the two pole square classes are necessarily different at a selected
prime. Requiring them to agree instead leads, for each fixed scale, to
a genus-four curve. Thus the established reduction identifies a precise
new realization problem; it does not settle that problem.
\end{abstract}

\section{Scope and parameter convention}
All equations in this paper have rational witnesses. For $a\ne0$,
$\sqcl a$ denotes its image in $\Q^\times/\Q^{\times2}$.
A \emph{retained parameter} is a rational coordinate not eliminated by
the finite-algebra step. An auxiliary square root is a finite algebraic
choice. It is not legitimate simply to omit the requirement that this
root be rational. Section~\ref{sec:packing} gives an explicit
one-coordinate polynomial test for a finite collection of such choices.
The existing finite-algebra constructions already use one separating
coordinate; it is shared by all their finite choices, rather than added
once per pole condition.

The conclusions below concern a source system in which pole-witness
coordinates occur only in their pole equations and stated nonvanishing
conditions. Further equations involving these coordinates would have to
be checked separately. Zariski density permits avoiding nonzero
polynomials; it does not permit imposing an arbitrary equation or
identifying arbitrary witness coordinates.

\section{The older pole system}
For $b\ne0$, put
\[
 B_b=17b,\qquad V_b=B_b(B_b-b^2)=17b^2(17-b).
\]
The older pole equations are
\begin{equation}\label{eq:old}
 p^2=B_b(1-e)(1+3e),\qquad
 B_bq^2=(s^2-eV_b)(s^2+3eV_b).
\end{equation}
On the usual conic chart,
\begin{equation}\label{eq:chart}
 e=\frac{\tau^2-B_b}{\tau^2+3B_b},\qquad
 p=-\frac{4B_b\tau}{\tau^2+3B_b}.
\end{equation}
The retained parameters are $(\tau,s)$, and the second equation specifies
a square root $q$. For positive $b$, the chart denominator never
vanishes at a rational $\tau$.

\begin{theorem}[Exact joint realization]\label{thm:old}
Let $r\ge1$, and let $b_0,\ldots,b_{r-1}$ belong to
$\Q\setminus\{0,17\}$. Suppose that, for every $i$,
\begin{equation}\label{eq:compatible17}
 b_i=z_i^2b_0,\qquad 17-b_i=c_i^2(17-b_0),
 \qquad z_i,c_i\in\Q^\times,
\end{equation}
where $z_0=c_0=1$. A solution $(e,p,s,q)$ of
\eqref{eq:old} at $b_0$ gives solutions at all the inputs by
\begin{equation}\label{eq:oldtransfer}
 e_i=e,\quad p_i=z_ip,\quad s_i=z_i^2c_i s,
 \quad q_i=z_i^3c_i^2q.
\end{equation}
On the conic charts, $\tau_i=z_i\tau$.
Consequently the joint pole conditions, restricted to
\eqref{eq:compatible17}, have a presentation with only the two retained
parameters $(\tau,s)$, together with finite quadratic choices.
The transformation preserves nonvanishing of $p,s,q$.
\end{theorem}
\begin{proof}
Write $B=B_{b_0}$ and $V=V_{b_0}$. The compatibility equations imply
\[
 B_{b_i}=z_i^2B,\qquad V_{b_i}=z_i^4c_i^2V.
\]
The first equation in \eqref{eq:old} therefore scales by $z_i^2$.
Moreover
\[
 s_i^2-e_iV_{b_i}=z_i^4c_i^2(s^2-eV),\qquad
 s_i^2+3e_iV_{b_i}=z_i^4c_i^2(s^2+3eV).
\]
Their product scales by $z_i^8c_i^4$, which is exactly the scaling of
$B_{b_i}q_i^2$. Substitution in \eqref{eq:chart} gives $e_i=e$ and
$p_i=z_ip$ when $\tau_i=z_i\tau$. All scalings are nonzero.
Conversely, the joint system includes its $b_0$ pole condition, so
existence of the joint system implies existence of that condition.
This proves both directions of the existential equivalence on the
specified compatibility locus.
\end{proof}

\paragraph{The retained realization input.}
An older prepared input has square reduced denominator and satisfies
\[
 0<b<17,\qquad b\in\Q_2^{\times2}\cap\Q_3^{\times2}\cap\Q_{17}^{\times2},
 \qquad v_2(b)\ge3,\quad v_3(b)\ge1,\quad v_{17}(b)\ge2.
\]
Write $V_b=\theta c^2$ with $\theta>0$ squarefree. The precise
arithmetic input used for density is existence of a positive $n$ and
infinite-order rational points on
\[
 E_n:y^2=x(x-n)(x+3n),\qquad E_{n\theta},
\]
both with abscissa in $B_b\Q^{\times2}$. This is the inherited
simultaneous realization from \cite{Density}; it is not proved anew here.

\begin{corollary}[Completeness and avoidance]\label{cor:oldcomplete}
Suppose $b_0$ is prepared for the inherited simultaneous pole-realization
input, in the precise sense of \cite{Density}. That input gives a
Zariski-dense set of witnesses in $(\tau,s)$ with $pqs\ne0$.
Under \eqref{eq:compatible17}, the same two parameters simultaneously
realize every $b_i$. Any finite collection of nonzero polynomials in
these two parameters may be required to be nonzero as well.
\end{corollary}
\begin{proof}
For clarity, the density assertion uses more than a black-box claim
that a single witness exists. Independently vary positive odd multiples
of the two supplied non-torsion points, and write their coordinates as
$(B u^2,y_1)$ and $(B v^2,y_2)$, where $B=B_{b_0}$ and
$V_{b_0}=\theta c^2$. Odd multiples preserve the abscissa square class,
as follows from the usual chord-and-tangent square-class identity.
The exact formulas
\[
 e=\frac{n}{Bu^2},\quad p=\frac{y_1}{Bu^3},\quad
 s=\frac{cv}{u},\quad q=\frac{c^2y_2}{B^2u^2v},\quad
 \tau=\frac{p}{e-1}
\]
give the base pole witness. Non-torsion gives $pqs\ne0$; positivity
of the first abscissa gives $0<e<1$. Varying the first point gives
infinitely many $\tau$, and, with it fixed, varying the second gives
infinitely many $s$. A polynomial vanishing on all such pairs is zero,
by first considering its coefficients as a polynomial in $s$, then
as polynomials in $\tau$. This proves density.
Apply Theorem~\ref{thm:old} to this single-input witness set.
The complement of a finite union of proper polynomial zero sets is a
nonempty Zariski-open subset of $\A^2$. It meets that set by the density
theorem of \cite{Density}. The relevant polynomials must remain nonzero
\emph{after} the substitutions \eqref{eq:oldtransfer}; this is a genuine
hypothesis, not a consequence of density of independently chosen
witnesses.
\end{proof}

\paragraph{The exact finite choices.}
After retaining $\tau,s$, use the radicands
\begin{equation}\label{eq:oldradicands}
 \frac{(s^2-eV_{b_0})(s^2+3eV_{b_0})}{B_{b_0}},\qquad
 \frac{b_i}{b_0},\qquad \frac{17-b_i}{17-b_0}\quad(1\le i<r).
\end{equation}
Their rational square roots give $q,z_i,c_i$. The first radicand may
be required nonzero by Corollary~\ref{cor:oldcomplete}; the others are
nonzero by the hypotheses. Thus $2r$ retained pole parameters are
replaced by two, at the cost of increasing the finite algebra.
For two pole conditions this saves two retained parameters. It is not
an assertion that the inputs in an existing complete construction
satisfy \eqref{eq:compatible17}.

\begin{proposition}[A necessary condition for a common conic coordinate]
If two solutions of the first equation in \eqref{eq:old} use the same
$e\notin\{1,-1/3\}$, then $b_1/b_0$ is a rational square.
\end{proposition}
\begin{proof}
Both $p_i$ are nonzero, and division of their equations gives
$b_1/b_0=(p_1/p_0)^2$. The excluded values permit $p_i=0$ and cannot
be included in this argument.
\end{proof}

\section{Congruent-number pole equations}
The pole system obtained from the Hexagon simultaneous-twist input is
\begin{equation}\label{eq:new}
 e^2+bp^2=1,\qquad uv=e(1-b),\qquad u^2-v^2=bq^2.
\end{equation}
Its local implication holds without any nonzero auxiliary restriction:
if $b\ne0$ and these equations have a rational solution, $v_\ell(b)$
is not negative and odd at an odd prime $\ell$ \cite{Congruent}.

\begin{theorem}[Transfer of an entire congruent-number witness]\label{thm:new}
Let $b_0,b_i\notin\{0,1\}$ and suppose
\begin{equation}\label{eq:compatible1}
 b_i=z_i^2b_0,\qquad 1-b_i=c_i^2(1-b_0),
 \qquad z_i c_i\ne0.
\end{equation}
Every witness $(e,p,u,v,q)$ at $b_0$ transfers to a witness at $b_i$ by
\begin{equation}\label{eq:newtransfer}
 e_i=e,\quad p_i=p/z_i,\quad u_i=c_i u,
 \quad v_i=c_i v,\quad q_i=c_i q/z_i.
\end{equation}
An arbitrary number of compatible inputs therefore needs only two
retained parameters, which may be taken to be $e,u$, and finite
quadratic choices on the nondegenerate locus.
\end{theorem}
\begin{proof}
The first residual is unchanged. The middle residual is multiplied by
$c_i^2$, and the last residual is multiplied by $c_i^2$. These identities
prove the assertion, including zero auxiliary values. On the locus
$b_0u\ne0$, the finite radicands for the base witness are
\[
 A=b_0(1-e^2),\qquad
 C=b_0\bigl(u^4-e^2(1-b_0)^2\bigr).
\]
Writing $a^2=A$ and $c^2=C$, one recovers
\[
 p=a/b_0,\qquad v=e(1-b_0)/u,\qquad q=c/(b_0u).
\]
The remaining choices are the square roots in \eqref{eq:compatible1}.
\end{proof}

\subsection{The arithmetic input and joint density}
Write $C_n:y^2=x^3-n^2x$. The precise arithmetic input used here is
Jaswal's Theorem~1.2 in Hexagon version~2 \cite{Jaswal}: if
$\beta>1$ and $\theta\ge1$ are coprime positive squarefree integers,
both $1\pmod8$, and each is a square modulo every prime dividing the
other, there is a positive squarefree $n$ for which $C_n$ and
$C_{n\theta}$ have infinite-order rational points with abscissa in
$\beta\Q^{\times2}$. This paper retains that statement as an input;
it is not a new independent audit of its proof.

If $0<b<1$ has square reduced denominator, is a square in $\Q_2$, and
$v_2(b)\ge3$, writing $b=\beta z^2$, $1-b=\theta c^2$ supplies exactly
these conditions when $\beta>1$. Coprimality and the mutual residue
conditions follow from the coprime numerators of $b$ and $1-b$ over a
square denominator. The square unit of $b$ at~2 gives
$\beta\equiv1\pmod8$; its even valuation is at least four, so
$1-b\equiv1\pmod {16}$ and $\theta\equiv1\pmod8$. These are also
proved in detail in \cite{Congruent}.

\begin{theorem}[Common first parameter under an explicit twist input]
\label{thm:commone}
Suppose $0<b_i<1$, $b_i=\beta z_i^2$ with a fixed positive squarefree
$\beta$, and $1-b_i=\theta_i c_i^2$.
Assume that a single positive $n$ is given such that $C_n$ and every
$C_{n\theta_i}$ have an infinite-order rational point with abscissa in
$\beta\Q^{\times2}$.
Then the systems \eqref{eq:new} for all $i$ have witnesses with a
\emph{common} $e$, with $0<e<1$ and $u_i p_i q_i\ne0$.
The realizable tuples $(e,u_0,\ldots,u_{r-1})$ are Zariski dense in
$\A^{r+1}$.
\end{theorem}
\begin{proof}
Choose independently positive odd multiples of the supplied points.
The square-class homomorphism on $C_n$, given away from two-torsion by
$(x,y)\mapsto\sqcl x$, is constant on odd multiples. For completeness,
the chord-and-tangent law gives $x(Q)x(R)x(Q+R)$ as a square when the
line has nonzero intercept, and
\[
 x(2Q)=\left(\frac{x(Q)^2+n^2}{2y(Q)}\right)^2.
\]
Adding $2Q$ successively proves the stated assertion for odd multiples
of a non-torsion point; none of the exceptional two-torsion cases can
occur in that sequence.

Write the selected first point as $(x_0,y_0)$ on $C_n$. For each $i$,
scale the selected point on $C_{n\theta_i}$ by $(x,y)\mapsto(c_i^2x,c_i^3y)$,
obtaining $(x_i,y_i)$ on $C_{n(1-b_i)}$. There are nonzero rational
$a_{0i},a_i$ with $x_0=b_i a_{0i}^2$, $x_i=b_i a_i^2$. Put
\[
 e=\frac n{x_0},\quad
 p_i=\frac{y_0}{b_i^2a_{0i}^3},\quad
 u_i=\frac{a_i}{a_{0i}},\quad
 v_i=\frac{e(1-b_i)}{u_i},\quad
 q_i=\frac{y_i}{b_i^2a_{0i}^3u_i^2}.
\]
The two curve equations prove \eqref{eq:new} by direct division.
The ordinates and abscissas are nonzero. Since $x_0>0$, its curve
 equation gives $x_0>n$; equality would be two-torsion. Thus $0<e<1$.

The first point has infinitely many distinct abscissas as its odd
multiple varies, hence gives infinitely many $e$. With it fixed,
each of the independently selected second points gives infinitely many
$u_i^2$, hence infinitely many $u_i$. A polynomial vanishing on all
these tuples first vanishes identically in the $u_i$ by the Cartesian
product of infinite sets, and then its coefficients vanish at infinitely
many $e$. This proves density.
\end{proof}

\begin{corollary}\label{cor:newdensity}
For prepared inputs satisfying \eqref{eq:compatible1} with a common
nontrivial class $\beta$, Jaswal's two-twist input suffices, without a
new multi-twist assumption. Reusing the same first and second points
and applying \eqref{eq:newtransfer} gives a Zariski-dense set of common
retained pairs $(e,u)$.
\end{corollary}
\begin{proof}
All complement classes are the same $\theta$. The construction in
Theorem~\ref{thm:commone}, with one independently varying first point and
one independently varying second point, gives the same two-variable
coefficient argument. The transfer formulas then give all inputs.
\end{proof}

\begin{remark}
For different complement classes, separate applications of a two-twist
theorem do not supply the \emph{same} $n$ required in
Theorem~\ref{thm:commone}. Its multi-twist hypothesis is explicit and is
not silently inferred from Jaswal's Theorem~1.2. If an input $b$ is a
rational square, it has the elementary witness $e=0$, $u=1$, $v=0$,
$p=q=1/\sqrt b$. If $1-b=c^2$ is a rational square, it has the witness
$e=1$, $p=q=0$, $u=v=c$. These exceptional strata need no retained
pole parameters, but neither gives dense witnesses on the generic
parameter plane.
\end{remark}

\section{Encoding the finite choices with one coordinate}\label{sec:packing}
The following elementary statement specifies the quantifier accounting.
It does not assert that a rational root of an arbitrary resultant
lifts to rational witnesses.

\begin{lemma}\label{lem:separator}
Let $a_1,\ldots,a_k$ be nonzero rationals, $k\ge1$, and put
\[
 L=4\left(1+\sum_i a_i^2\right)
       \left(1+\sum_i a_i^{-2}\right),\qquad
 \Omega=1+\sum_{i=1}^k L^i Z_i
\]
in $A=\Q[Z_1,\ldots,Z_k]/(Z_i^2-a_i)$.
All $2^k$ geometric values of $\Omega$ are distinct and nonzero.
The characteristic polynomial of multiplication by $\Omega$ has a
rational root if and only if every $a_i$ is a rational square.
\end{lemma}
\begin{proof}
Let $m=\min_i\sqrt{|a_i|}$ and $M=\max_i\sqrt{|a_i|}$. The two
factors defining $L$ show
\[
 L\ge4,\qquad L\ge4M/m,\qquad L\ge4/m.
\]
For two distinct sign choices, let $j$ be their largest differing
index. Its contribution has modulus $2L^jm$. The smaller contributions
have total modulus at most
$2M\sum_{i<j}L^i<2ML^j/(L-1)\le 2mL^j/3$.
They cannot cancel it. The same estimate, together with
$1/L^k\le1/L\le m/4$, proves that the largest term of any value of
$\Omega$ cannot be canceled by the remaining terms and the constant~1.
These modulus arguments include negative radicands.

If a geometric value is rational, every automorphism of $\overline\Q/\Q$
fixes it. Distinctness forces that automorphism to fix each selected
$Z_i$, so all the square roots are rational. Conversely, if all roots
are rational, all geometric values are rational. The algebra is finite
\'etale because each $a_i\ne0$, even when their square classes are
linearly dependent.
\end{proof}

\begin{proposition}\label{prop:onecoordinate}
Let the $a_i$ be nonzero rational functions of retained coordinates
$\boldsymbol t$, and let $g(\boldsymbol t)\ne0$ describe any additional
polynomial nonvanishing conditions on the intended open locus.
There is an integer polynomial $P(\boldsymbol t,Y)$, with only one new
rational coordinate $Y$, whose rational projection is precisely the
locus where all input denominators, all $a_i$, and $g$ are nonzero and
all $a_i$ are rational squares.
\end{proposition}
\begin{proof}
Form the multiplication matrix $M_\Omega$ in the monomial basis
$Z_1^{\epsilon_1}\cdots Z_k^{\epsilon_k}$. Choose a nonzero integer
polynomial $D$ which clears its denominators and is a multiple of the
numerators and denominators of all $a_i$, of all input denominators,
and of $g$. Increasing powers of those factors if needed achieves both
requirements. Put $C=D M_\Omega$, an integer-polynomial matrix, and
\[
                         P=\det(I-YDC).
\]
If $D=0$, the polynomial has value 1, so no unwanted boundary solution
is introduced. If $D\ne0$, its geometric zeros have
$Y=(D^2\omega)^{-1}$, where $\omega$ is a geometric value of $\Omega$.
Lemma~\ref{lem:separator} says exactly when one such $Y$ is rational.
\end{proof}

For a stand-alone block of $r$ older pole conditions, independent
charts retain $2r$ parameters and one common coordinate for their square
roots. Theorem~\ref{thm:old} instead retains two parameters and one
common coordinate for \eqref{eq:oldradicands}, on the compatibility
locus. Thus the auxiliary count is $2r+1$ versus $3$, not $3r$ versus
$3$. Inside a larger finite-algebra construction the existing separating
coordinate is reused; the savings in retained coordinates are
$2r-2$. The other finite algebra and its guards must still be included.
No quartic-degree claim follows from this general polynomial encoding.

\section{Why the present prime construction cannot supply the match}
The six-contact conditional construction uses
\begin{equation}\label{eq:F}
 \Theta=\left\{0,-\tfrac12,-\tfrac14,-\tfrac25,-1,-\tfrac13\right\},
 \quad F(t)=\prod_{\theta\in\Theta}(t-\theta),\quad
 t_0=\lambda,\quad t_1=a\lambda.
\end{equation}
Here $a=\beta^d+T^d>4$, with the notation and hypotheses of
\cite{Input13,UnitNorm}. The nonzero pole inputs have the form
\begin{equation}\label{eq:actualclass}
 b_j=F(t_j)\left(\frac{\N(W_j)}{2(t_j-1)^4(W_j)_0}\right)^2.
\end{equation}
Thus $\sqcl {b_j}=\sqcl {F(t_j)}$. This statement is about the literal
source equations, not a numerical approximation to them.

\begin{proposition}[Incompatibility of the selected prime classes]
\label{prop:fresh}
In the simultaneous-prime completeness construction for
\eqref{eq:F}, suppose each distinct contact coefficient has its own
selected prime, with valuation one in its incident contact numerator,
and with no incidence in any other distinct coefficient. These are the
retained fresh-prime properties of that construction. Then
$b_1/b_0$ in \eqref{eq:actualclass} is not a rational square.
\end{proposition}
\begin{proof}
The nonzero numbers in $\Theta$ have absolute values between $1/4$ and~1.
If $a>4$, the nonzero elements of $\Theta/a$ have absolute value less
than $1/4$. Hence the contact coefficient sets $\Theta$ and $\Theta/a$
intersect only at~0. Choose a selected prime belonging to a nonzero
contact of the first detector. It has valuation one in $F(t_0)$ and
valuation zero in $F(t_1)$; the fixed scale and denominator primes are
excluded in the retained prime construction. The ratio in
\eqref{eq:actualclass} therefore has odd valuation at this prime.
\end{proof}

This proves an incompatibility of the \emph{existing completeness
witnesses} with both proposed first-class matching conditions. It does
not prove that every possible witness has different classes, or that
thirteen variables is a lower bound. A new complete family, not a
renaming of the present parameters, is needed.

\begin{proposition}[The fixed-scale curve]\label{prop:genus}
Fix a rational $a>4$. Write $F(t)=tG(t)$, with $G$ the monic degree-five
polynomial defined by \eqref{eq:F}. For nonzero $\lambda$ away from
all contacts,
\[
 \sqcl {F(\lambda)}=\sqcl {F(a\lambda)}
 \quad\Longleftrightarrow\quad
 \exists y\in\Q:\ y^2=aG(\lambda)G(a\lambda).
\]
The smooth projective model of this curve has genus four. In particular,
for each fixed rational $a>4$ it has only finitely many rational
$\lambda$, by Faltings' theorem.
\end{proposition}
\begin{proof}
Two nonzero rationals have the same square class exactly when their
product is a square, and
$F(\lambda)F(a\lambda)=a\lambda^2G(\lambda)G(a\lambda)$.
The ten roots are the nonzero contacts and their quotients by~$a$.
They are distinct by the argument in Proposition~\ref{prop:fresh}.
A squarefree polynomial of degree ten defines a double cover of
$\mathbf P^1$ branched at ten points and unramified at infinity.
Riemann--Hurwitz gives $2g-2=-4+10$, so $g=4$.
Faltings' theorem applies to this smooth projective curve over $\Q$.
\end{proof}

\begin{remark}
The assertion is pointwise in $a$, not a uniform bound on rational
points as $a$ varies. It does not rule out one suitable finite witness
for every required input. It does rule out proving the desired match by
an arbitrarily-large-slope argument for a fixed scale. The points at
infinity do not give finite witnesses to the source equations.
\end{remark}

\subsection{The ten-contact and three-pole variants}
A similar calculation applies to the twenty-five-variable family.
At fixed values of its other retained coordinates, write
\[
 \Xi=\{0,-1/2,-1,1/2399,-32/407\},\qquad
 \Xi_j=(4j+2+\Xi)\cup(4j+4+\Xi),
\]
\[
 H_j(\lambda)=\prod_{\xi\in\Xi_j}(a_j\lambda-\xi),\qquad a_0=1.
\]
Each $\Xi_j$ has ten distinct nonzero elements. The literal pole
input is $\rho H_j$ times a nonzero rational square, so equality of
first square classes is equivalent to equality of the classes of
$H_j$. The common factor $\rho$ cancels.

\begin{proposition}\label{prop:multigenus}
Let $r\ge2$, and let $H_0,\ldots,H_{r-1}$ be squarefree polynomials
of the same positive even degree $m$, with disjoint root sets over
$\overline\Q$. The smooth projective model of the simultaneous
square-class correspondence
\[
 y_i^2=H_0(\lambda)H_i(\lambda)\quad(1\le i<r)
\]
is geometrically connected and has genus
\[
                         g=1+2^{r-3}(rm-4).
\]
For the ten-contact family, outside the finite set of contact-ratio
collisions this gives genus $9$ for two poles and genus $27$ for three.
\end{proposition}
\begin{proof}
The $r-1$ functions $H_0H_i$ are independent modulo squares in
$\overline\Q(\lambda)^\times$: a root unique to $H_i$ detects its
exponent in any putative square product. The function-field extension
therefore has degree $d=2^{r-1}$ and is geometrically connected.
At each of the $rm$ distinct roots, the nonzero parity vector gives
inertia of order two. Thus each branch point contributes $d/2$ to
Riemann--Hurwitz. There is no ramification at infinity because each
$H_0H_i$ has even degree. Hence
\[
             2g-2=-2d+rm\,d/2,
\]
which is the displayed formula. In the ten-contact application,
root sets meet exactly when $a_j/a_i=\xi_j/\xi_i$ for some contacts;
this is a finite exceptional collection of ratios. Taking $m=10$
and $r=2,3$ gives the asserted genera.
\end{proof}

The statement freezes the scales and all other retained variables
before varying the slope. It does not assert finiteness in the whole
parameter space. The additional complement-class requirements in
Theorem~\ref{thm:old} only restrict this correspondence further.
As in Proposition~\ref{prop:genus}, the finiteness conclusion is an
obstruction to retaining the old unlimited-slope proof, not a proof
that no finite choice can work.

\section{An exact example and verification scope}
For the newer system take
\[
 b_0=\frac{4880}{9801},\quad e=\frac{41}{81},\quad
 p=\frac{11}{9},\quad u=\frac{41}{81},\quad
 v=\frac{4921}{9801},\quad q=\frac1{11}.
\]
These are an exact nondegenerate rational solution of \eqref{eq:new}.
Let $z=-1220/6141$ and $c=-8581/6141$. Then
\[
 b_0z^2+(1-b_0)c^2=1,\qquad
 b_1=b_0z^2=\frac{7263392000}{369614145681}.
\]
Both inputs are prepared, and \eqref{eq:newtransfer} gives their shared
witness. Their common squarefree representatives are
$\beta=305$ and $\theta=4921=7\cdot19\cdot37$.
The example is not evidence for completeness at arbitrary input
classes; that is the separate theorem and hypothesis discussion above.

The reproducible checker in
\nolinkurl{computation/pole_parameters/joint_realization_v29/check.py}
verifies the two transfer identities in polynomial rings, the elliptic
scaling formulas modulo the curve equations, the displayed rational
example, finite-algebra separation examples, deliberate corruption
controls, and the fixed-scale root computation. It authenticates the
retained source files by hash. These finite checks do not prove
Jaswal's theorem, the inherited pole-realization theorem, Faltings'
theorem, or completeness of a smaller global construction.

\section{Conclusions for the current constructions}
Theorems~\ref{thm:old} and~\ref{thm:new} establish genuine parameter
reductions for compatible pole inputs. They do not establish the
compatibility of the two inputs in the current twenty-five- or
thirteen-variable definitions. Proposition~\ref{prop:fresh} proves
that the present thirteen-variable completeness witnesses fail even
the first matching condition. Proposition~\ref{prop:genus} identifies
a fixed-scale curve governing any replacement with matching first
classes. No result table is changed: the currently established general
counts remain twenty-five and thirteen, with their previously stated
arithmetic hypotheses. The missing assertion is existence of a
suitable complete restricted witness family, not the algebraic transfer
or the quantifier accounting.

\begin{thebibliography}{9}
\bibitem{Density} \emph{Zariski density of the prepared pole witnesses},
project manuscript, 9 November 2026,
\nolinkurl{papers/pole_parameters/pole_witness_density_20261109_v14.tex}.
\bibitem{Congruent} \emph{Pole equations from simultaneous congruent-number twists},
project manuscript, 14 November 2026,
\nolinkurl{papers/pole_parameters/congruent_twist_poles_20261114_v27.tex}.
\bibitem{Jaswal} M.~Jaswal, \emph{A proof that $\Z$ is Diophantine in $\Q$
through congruent-number twists}, Hexagon 2611.00331v2,
12 November 2026, Theorem~1.2.
\url{https://hexagonmath.org/2611.00331v2}.
\bibitem{Input13} \emph{Thirteen Variables from Direct Elliptic Index Comparison},
project source manuscript,
\nolinkurl{papers/conditional_few_variables/direct_input13_20261110_v22.tex}.
\bibitem{UnitNorm} \emph{A single norm determinant for the 25- and 13-variable constructions},
project manuscript, 14 November 2026,
\nolinkurl{papers/finite_algebra/unit_norm_20261114_v28.tex}.
\bibitem{Faltings} G.~Faltings, \emph{Endlichkeitss\"atze f\"ur abelsche
Variet\"aten \"uber Zahlk\"orpern}, Invent. Math. \textbf{73} (1983),
349--366, Satz~7, p.~365; erratum, \textbf{75} (1984), 381.
\end{thebibliography}
\end{document}
